\documentclass[11pt]{article}

\usepackage[T1]{fontenc}
\usepackage{amsmath,amssymb,amsthm}
\usepackage[margin=1in]{geometry}
\usepackage{hyperref}

\newtheorem{proposition}{Proposition}

\title{Regime-Aware Spectral Risk Estimation\\
for Cost-Constrained Statistical Arbitrage}
\author{Your Name}
\date{\today}

\begin{document}
\maketitle

\begin{abstract}
This paper studies whether random matrix theory (RMT) covariance regularization
improves the risk allocation of a cost-constrained statistical arbitrage strategy.
The proposed strategy separates the return signal from the risk model: a
pre-specified residual mean-reversion signal determines expected returns, while
spectral covariance estimates determine portfolio risk. We propose a
regime-aware estimator that blends long- and short-window cleaned covariance
matrices using a rolling measure of correlation drift. The estimator's
mathematical validity and empirical value are separate questions. We establish
positive semidefiniteness of the blend and specify a walk-forward design to test
whether it improves out-of-sample risk and net performance. No empirical result
is claimed in this draft.
\end{abstract}

\section{Research question and contribution}

The central question is whether adapting the covariance risk model to detected
changes in correlation structure improves a statistical arbitrage portfolio after
transaction costs.

The proposed extension is a regime-aware blend of long- and short-window RMT
covariance estimates. The blend is designed to retain the stability of a longer
estimation window while responding to changes indicated by a shorter window.
Its usefulness is an empirical hypothesis, not an assumed result. The literature
review must establish whether this exact construction and evaluation have already
been studied.

\section{Signal and portfolio model}

Let $r_t \in \mathbb{R}^N$ denote the vector of asset returns at time $t$.
A separate signal model produces a forecast $\widehat{\mu}_t \in \mathbb{R}^N$.
For example, asset returns may first be decomposed into factor and residual
components:
\[
    r_t = B f_t + \varepsilon_t,
\]
where $f_t$ is a vector of common factors, $B$ is a factor loading matrix, and
$\varepsilon_t$ contains residual returns.

A possible mean-reversion signal uses standardized residuals:
\[
    z_{i,t} = \frac{\widehat{\varepsilon}_{i,t}}{\widehat{s}_{i,t}},
    \qquad
    \widehat{\mu}_{i,t} = -\theta_i z_{i,t}.
\]
The factor model, residual scale, and parameters $\theta_i$ must be estimated
using information available at time $t$ only. This signal specification is a
starting point; its predictive value must be tested independently of the
covariance estimator.

Given a covariance estimate $\widehat{\Sigma}_t$, portfolio weights are chosen
by solving
\[
\widehat{w}_t =
\arg\max_{w \in \mathcal{W}}
\left\{
    \widehat{\mu}_t^\top w
    - \frac{\lambda}{2} w^\top \widehat{\Sigma}_t w
    - c \lVert w-w_{t-1}\rVert_1
\right\}.
\]
Here $\lambda > 0$ controls risk aversion, $c \geq 0$ represents proportional
trading costs, and $\mathcal{W}$ contains leverage, position, and exposure
constraints. These choices must be fixed or calibrated using training data, not
the final test period.

\section{Spectral covariance estimation}

Let $X \in \mathbb{R}^{T \times N}$ be a rolling matrix of demeaned,
standardized returns, with $T$ observations and $N$ assets. The sample
correlation matrix is
\[
    C = \frac{1}{T}X^\top X.
\]
Write its eigendecomposition as
\[
    C = U \operatorname{diag}(\lambda_1,\ldots,\lambda_N) U^\top.
\]

Under an idealized null model of independent standardized observations, the
Marchenko--Pastur distribution describes the limiting eigenvalue spectrum as
$N,T \to \infty$ with $q=N/T$ fixed. For unit population variance, the bulk
support is
\[
    \lambda_- = (1-\sqrt{q})^2,
    \qquad
    \lambda_+ = (1+\sqrt{q})^2.
\]
This null model is a reference, not a complete model of financial returns.
Serial dependence, heavy tails, and changing correlations can affect the
spectrum and must be considered in robustness analysis.

As a baseline cleaning rule, eigenvalues inside the estimated noise bulk may be
replaced by their average, while selected eigenvalues above the bulk are retained.
The cleaned correlation matrix is reconstructed from the original eigenvectors
and transformed eigenvalues. Marginal volatilities are then restored to obtain a
covariance estimate.

\section{Proposed regime-aware extension}

At each time $t$, construct cleaned correlation estimates from long and short
rolling windows, denoted by $\widetilde{C}^{(L)}_t$ and
$\widetilde{C}^{(S)}_t$. Define their relative drift as
\[
    d_t =
    \frac{
        \left\lVert \widetilde{C}^{(S)}_t -
        \widetilde{C}^{(L)}_t \right\rVert_F
    }{
        \left\lVert \widetilde{C}^{(L)}_t \right\rVert_F + \epsilon
    },
\]
where $\epsilon > 0$ prevents division by zero. For a scale parameter
$\tau > 0$, set
\[
    \alpha_t = \frac{d_t}{d_t+\tau},
    \qquad
    \widetilde{C}^{(A)}_t =
    (1-\alpha_t)\widetilde{C}^{(L)}_t
    + \alpha_t \widetilde{C}^{(S)}_t.
\]
The parameter $\tau$, window lengths, and all preprocessing choices must be
selected using training data only. The study will compare this rule with
fixed-window RMT cleaning and established shrinkage baselines.

\begin{proposition}
If $\widetilde{C}^{(L)}_t$ and $\widetilde{C}^{(S)}_t$ are positive
semidefinite and $0 \leq \alpha_t \leq 1$, then
$\widetilde{C}^{(A)}_t$ is positive semidefinite.
\end{proposition}

\begin{proof}
For any vector $x \in \mathbb{R}^N$,
\[
x^\top \widetilde{C}^{(A)}_t x
=
(1-\alpha_t)x^\top \widetilde{C}^{(L)}_t x
+\alpha_t x^\top \widetilde{C}^{(S)}_t x.
\]
Both quadratic forms on the right are nonnegative, and both coefficients are
nonnegative. Therefore
$x^\top \widetilde{C}^{(A)}_t x \geq 0$ for every $x$, which proves the claim.
\end{proof}

This proposition establishes mathematical validity of the blend. It does not
establish that the estimator is more accurate or profitable.

\section{Empirical hypotheses}

The walk-forward study will test the following hypotheses:

\begin{itemize}
    \item[H1.] The regime-aware estimate reduces realized out-of-sample portfolio
    variance relative to the sample covariance baseline.
    \item[H2.] Any risk improvement persists after turnover and transaction costs.
    \item[H3.] The estimated benefit is robust to reasonable changes in window
    lengths, market periods, and cost assumptions.
\end{itemize}

These are hypotheses to test, not conclusions.

\section{Evaluation design}

All preprocessing, universe membership, parameter selection, signal estimation,
and portfolio construction must use only information available at each decision
date. Use nested walk-forward evaluation: an inner training period selects
parameters, and a later untouched period evaluates them.

Baselines should include equal weighting, sample covariance, a standard
linear-shrinkage estimator, fixed-window RMT cleaning, and the proposed
regime-aware blend. Report realized volatility, drawdown, turnover, gross and
net performance, and uncertainty across time periods. Include transaction costs
and short-borrow assumptions where applicable.

The final study should disclose data sources, survivorship and look-ahead
controls, missing-data rules, rebalance timing, and all parameter choices.

\section{Results}

\textit{To be completed after the empirical protocol is fixed and the experiments
are run. Do not insert performance claims before then.}

\section{Limitations}

The Marchenko--Pastur benchmark assumes an idealized sampling model. Financial
returns can violate its assumptions through serial dependence, heavy tails,
heteroskedasticity, and structural breaks. The proposed drift statistic can also
react to estimation noise. These issues motivate robustness checks; they do not
guarantee that the proposed estimator will outperform its baselines.

\begin{thebibliography}{9}

\bibitem{marchenko-pastur}
V. A. Mar\v{c}enko and L. A. Pastur.
Distribution of eigenvalues for some sets of random matrices.
\textit{Mathematics of the USSR-Sbornik}, 1(4):457--483, 1967.

\bibitem{laloux}
L. Laloux, P. Cizeau, J.-P. Bouchaud, and M. Potters.
Noise dressing of financial correlation matrices.
\textit{Physical Review Letters}, 83:1467--1470, 1999.

\bibitem{markowitz}
H. Markowitz.
Portfolio selection.
\textit{The Journal of Finance}, 7(1):77--91, 1952.

\bibitem{ledoit-wolf}
O. Ledoit and M. Wolf.
Nonlinear shrinkage of the covariance matrix for portfolio selection:
Markowitz meets Goldilocks.
\textit{The Review of Financial Studies}, 30(12):4349--4388, 2017.

\end{thebibliography}

\end{document}